\einheitenkopf{Einheit 5 von 5 · Anwendungen}
\zweigzeile{Hier lernst du, Tarife und andere Sachsituationen mit linearen Funktionen zu beschreiben, zu berechnen und zu vergleichen · neu in diesem Jahr · P10 oft · baut auf: Gleichung, Graph und Schnittpunkt (Einheit 2 bis 4), Dreisatz}
\erg{41}{Anfangswert – Änderung je Einheit: a) 8 € – 9 Cent je Minute \quad b) 500 Liter – 20 Liter weniger je Stunde \quad c) 29 € – 15 € je Monat \quad d) 3,90 € – 2,10 € je km \quad e) 50 cm – 3 cm je Monat}
\erg{42}{a) $y = 0{,}3x + 6$ \quad b) $y = 4x + 12$ \quad c) $y = 0{,}3x + 9$ (30 Cent = 0,30 €)}
\erg{43}{a) $y = 3x + 5$ \quad b) $y = -25x + 800$ \quad c) $y = 0{,}08x + 2$}
\erg{44}{a) $40 + 6 \cdot 15 = 130$, also 130 € \quad b) $60 - 7{,}5 \cdot 6 = 15$, also 15 m³ \quad c) je Woche 4 cm weniger: $y = -4x + 120$, $120 - 4 \cdot 10 = 80$, also 80 cm \quad d) $45x + 350 > 1000$ ergibt $x > 14{,}\overline{4}$; nach 14 Monaten 980 €, nach 15 Monaten 1025 €: nach 15 Monaten}
\erg{45}{a) A: $6 + 0{,}08 \cdot 100 = 14$ €, B: $0{,}15 \cdot 100 = 15$ €; A ist günstiger. \quad b) C: $12 + 0{,}20 \cdot 70 = 26$ €, D: $4 + 0{,}10 \cdot 120 = 16$ €; D ist günstiger.}
\erg{46}{a) A: $y = 2{,}5x + 8$, B: $y = 4{,}5x$ \quad b) $2{,}5x + 8 = 4{,}5x$, $x = 4$; bei 4 Stunden kosten beide 18 €, ab mehr als 4 Stunden ist A günstiger. \quad c) ja: A kostet 33 €, B 45 €; nach dem Schnittpunkt steigt B je Stunde um 2 € mehr als A.}
\erg{47}{a) Tim hat Anfangswert und Änderung vertauscht. Richtig: $y = 35x + 20$, $y = 35 \cdot 6 + 20 = 230$, also 230 €. \quad b) $y = 0{,}6x + 3{,}5$; $0{,}6 \cdot 8 + 3{,}5 = 8{,}30$, also 8,30 €}
\erg{48}{a) A: III, B: I, C: II \quad b) B hat keine Grundgebühr, der Graph beginnt im Ursprung; das ist nur bei I so. \quad c) B (3 €; A 4 €, C 6 €)}
\erg{49}{Beispiele: a) Fahrradverleih: 10 € Grundgebühr ($n$), 2,50 € je Stunde ($m$); $x$ Stunden, $y$ Preis in € \quad b) Wassertank mit 300 Litern ($n$), der je Tag 15 Liter verliert ($m$); $x$ Tage, $y$ Liter \quad c) Äpfel zu 0,12 € je Stück ohne Grundpreis ($n = 0$); $x$ Stück, $y$ Preis in €}
